\documentclass{article}
\usepackage{graphicx} % Required for inserting images
\usepackage[T1]{fontenc}
\usepackage[utf8]{inputenc}
\usepackage[polish]{babel}
\usepackage{lmodern}
\usepackage[hidelinks]{hyperref} 
\usepackage{booktabs}
\usepackage{csvsimple}
\usepackage[table]{xcolor}
\usepackage{array}
\usepackage{colortbl}
\usepackage{amsmath}
\usepackage{geometry}
\usepackage{pdfpages}
\usepackage{ragged2e}

\begin{document}

\begin{titlepage}
\newgeometry{left=1.6in, right=2in}
\vspace*{1.5cm}

\noindent  \textcolor{gray}{\large Marszałek, Marynowski, Winowiecki} \\
\vspace{1cm}

\noindent \textbf{\Large Część laboratoryjna PAiR - Ćwiczenie 1} \\
\vspace{0.5cm}

\noindent {\large Podstawowe człony układów regulacji automatycznej} \\


Listopad 2026 \\

\vspace{0.4cm}
\noindent Politechnika Gdańska \\
Wydział Elektroniki, Telekomunikacji i Informatyki \\
Katedra Inżynierii Biomedycznej \\

\begin{figure}[h]
    \includegraphics[width=0.33\textwidth]{logol.jpg}
\end{figure}
\end{titlepage}




\section*{Aparatura użyta do wykonania zadania}

\section*{Sprawozdanie właściwe}

\section{Badanie modeli fizycznych}
    
    \subsection{Człon proporcjonalny $G(s)=k_0$}
        \includegraphics[scale=0.5]{1real.pdf}

    $k_0=\frac{y_{ust}}{A}=\frac{8}{4,4}\approx 1,82$
    
    \subsection{Człon inercyjny pierwszego rzędu $G(s)=\frac k{1+sT_1}$}
        \includegraphics[scale=0.5]{2real.pdf}

    $k=\frac{y_{ust}}{A}=\frac{8,8}{4,4}=2$
    
    $T_1=(2,7 - 2)\times200ms=0,14s$
    
    \subsection{Człon całkujący idealny i rzeczywisty}
        
        \subsubsection{Człon całkujący idealny: $G(s)=\frac k{s}$}
            \includegraphics[scale=0.5]{3ideal.pdf}

        $k=\frac{y_{ust}}{A\times \Delta t}=\frac{10,8}{4,4\times0,4}\approx6,136$
        
        \subsubsection{Człon całkujący rzeczywisty: $G(s)= \frac k {s(1+sT_1)}$}
            \includegraphics[scale=0.5]{3real.pdf}

        $k=\frac{y_{ust}}{A\times \Delta t}=\frac{10,8}{4,4\times0,46}\approx5,336$

        $T_1=(2,6 - 2)\times200ms=0,12s$
    
    \subsection{Korektor szeregowy z opóźnieniem fazy $G(s) = k\frac{1+sT_1}{1+sT_2}, T_1<T_2$}
        \includegraphics[scale=0.5]{4real.pdf}

        $k=\frac{y_{ust}}{A}=\frac{10,8}{4,4}\approx1,82$

        $T_1=0,1s$

        $T_2=0,2s$
    
    \subsection{Korektor szeregowy z wyprzedzeniem fazy $G(s) = k\frac{1+sT_1}{1+sT_2}, T_1<T_2$}
        \includegraphics[scale=0.5]{5real.pdf}

        $k=\frac{y_{ust}}{A}=\frac{4}{4,4}\approx0,91$

        $T_1=0,4s$

        $T_2=0,2s$
    
    \subsection{Człon oscylacyjny drugiego rzędu $G(s)=\frac{k\omega_0^2}{s^2+2D\omega_0s+\omega_0^2}, 0<D<1$}
        \includegraphics[scale=0.5]{6real.pdf}

        $k=\frac{y_{ust}}{A}=\frac{5,5}{4,4}=1,25$

        $\omega_t=6,56 \frac{1}{s}$

        $D=0,416$

        $\omega_0=7,21\frac{1}{s}$
    
    \subsection{Człon inercyjny drugiego rzędu $G(s)=\frac{k}{(1+sT_1)(1+sT_2)}$}
        \includegraphics[scale=0.5]{7real.pdf}

        $k=\frac{y_{ust}}{A}=\frac{8}{4,4}\approx1,82$

        $T_1=0,63s$

        $T_2=0,24s$
        
        %tutaj mamy dać jeszcze opis
        
\section{Badanie własnych układów}

\section{Badanie liniowych układów automatyki\newline na drodze modelowania (MATLAB)}
    
    \subsection{Badanie członu inercyjnego pierwszego rzędu w oparciu o zadane parametry}
    
        \subsubsection{$k=1, T_1=0.7$}
            \includegraphics[scale=0.5]{1_1-1.pdf}
                $y_{ust}=1\\
                k=\frac{y_{ust}}{A}=\frac{1}{1}\\
                y(T)=0.632\cdot1=0.632\\
                T\approx0.7s\\$
            \includegraphics[scale=0.5]{1_1-2.pdf}
                $g(0^+)=\frac{k}{T_1}=\frac{1}{0.7}\approx1.43\\
                g(T)=0.368\cdot1.43=0.526\\
                T\approx0.7s\\$
            \includegraphics[scale=0.5]{1_1-3.pdf}
                $k=|G(j0)|=1\\
                |G(j\omega_1)|=0.707\cdot1=0.707\\
                \omega_1\approx1.43\text{ rad/s}\\
                T=\frac{1}{\omega_1}\approx\frac{1}{1.43}\approx0.7s\\$
            \includegraphics[scale=0.5]{1_1-4.pdf}
                $\varphi(\omega_2)=-\frac{\pi}{4}=-45^\circ\\
                \omega_2\approx1.43\text{ rad/s}\\
                T=\frac{1}{\omega_2}\approx\frac{1}{1.43}\approx0.7s\\$
        
        \subsubsection{$k=1, T_1=1.2$}
            \includegraphics[scale=0.5]{1_1-5.pdf}
                $y_{ust}=1\\
                k=\frac{y_{ust}}{A}=\frac{1}{1}\\
                y(T)=0.632\cdot1=0.632\\
                T\approx1.2s\\$  
            \includegraphics[scale=0.5]{1_1-6.pdf}
                $g(0^+)=\frac{k}{T_1}=\frac{1}{1.2}\approx0.833\\
                g(T)=0.368\cdot0.833=0.307\\
                T\approx1.2s\\$
            \includegraphics[scale=0.5]{1_1-7.pdf}
                $k=|G(j0)|=1\\
                |G(j\omega_1)|=0.707\cdot1=0.707\\
                \omega_1\approx0.83\text{ rad/s}\\
                T=\frac{1}{\omega_1}\approx\frac{1}{0.833}\approx1.2s\\$
            \includegraphics[scale=0.5]{1_1-8.pdf}
                $\varphi(\omega_2)=-\frac{\pi}{4}=-45^\circ\\
                \omega_2\approx0.83\text{ rad/s}\\
                T=\frac{1}{\omega_2}\approx\frac{1}{0.833}\approx1.2s\\$
        
    \subsection{Badanie członu oscylacyjnego drugiego rzędu w oparciu o zadane parametry}
        
        \subsubsection{$k=1, D=0.3, \frac{T_0}{2\pi}=0.5$}
            \includegraphics[scale=0.5]{2_1-1.pdf}
    
            \includegraphics[scale=0.5]{2_1-2.pdf}
    
            \includegraphics[scale=0.5]{2_1-3.pdf}
    
        
        \subsubsection{$k=1, D=0.5, \frac{T_0}{2\pi}=0.5$}
            \includegraphics[scale=0.5]{2_1-4.pdf}
    
            \includegraphics[scale=0.5]{2_1-5.pdf}
    
            \includegraphics[scale=0.5]{2_1-6.pdf}
    
        
        \subsubsection{$k=1, D=0.3, \frac{T_0}{2\pi}=1.5$}
            \includegraphics[scale=0.5]{2_1-7.pdf}
    
            \includegraphics[scale=0.5]{2_1-8.pdf}
    
            \includegraphics[scale=0.5]{2_1-9.pdf}
    
        
\section{Opracowanie wyników}

\section*{Wnioski z pracy}

\end{document}
